\section*{ВВЕДЕНИЕ}
\addcontentsline{toc}{section}{ВВЕДЕНИЕ}

В условиях устойчивой цифровизации экономики и социальной сферы все больше процессов управления взаимодействием с клиентами переводится в гибридный и полностью цифровой формат. Для учебных проектов по направлению разработки программного обеспечения это означает смещение фокуса от простых демонстрационных приложений к системам, которые отражают реальный жизненный цикл продукта: учет требований, архитектурное проектирование, разработка пользовательского интерфейса, интеграция с серверной частью, контроль качества и подготовка к развертыванию. В этой логике система \textit{Gamified Task CRM (Flutter Web)} рассматривается как прикладной программный продукт, объединяющий функции управления задачами, базовые CRM-процессы, модуль отчетности и вспомогательные интеллектуальные сервисы \cite{repoReadme,docsIndex}.

Актуальность темы курсового проекта определяется сразу несколькими факторами. Во-первых, бизнес-пользователи ожидают от CRM-решений не только хранения карточек клиентов, но и связности всех операционных сущностей: задач, этапов работы, аналитики и персональных метрик. Во-вторых, активно растет спрос на кроссплатформенные технологии, позволяющие быстро выпускать единый продукт для web-окружения без отдельной реализации на нескольких стековых решениях \cite{flutterDocs,dartDocs}. В-третьих, в образовательном процессе становится критичным умение строить приложение на современных практиках: декларативный UI, управление состоянием, строгая маршрутизация, безопасная работа с конфигурацией и воспроизводимый тестовый контур \cite{riverpodDocs,goRouterDocs,owaspM10}. Таким образом, выбранная тема позволяет связать теоретические положения дисциплины с практически значимой реализацией.

Дополнительный аргумент в пользу выбранной темы состоит в том, что разрабатываемая система находится на стыке нескольких предметных областей: проектирование интерфейсов, управление данными, прикладная аналитика, инженерия качества и основы информационной безопасности. Именно такие комплексные проекты позволяют продемонстрировать уровень профессиональной подготовки студента, поскольку требуют не фрагментарных знаний, а целостного подхода к разработке программного обеспечения \cite{sommerville,cleanArchitecture}.

\textbf{Целью} курсового проекта является разработка и обоснование архитектурно-целостного web-приложения \textit{Gamified Task CRM} для управления задачами и клиентскими взаимодействиями с использованием современного Flutter-стека и backend-first интеграции с Supabase.

Для достижения поставленной цели в работе решаются следующие задачи:
\begin{enumerate}
  \item Проанализировать теоретические и технологические основы построения CRM-ориентированных систем управления задачами, включая требования к архитектуре, данным и пользовательскому опыту.
  \item Спроектировать логическую структуру программного решения: модули интерфейса, маршрутизацию, модель состояния, состав ключевых сущностей и взаимодействие с backend-сервисами.
  \item Реализовать практическую часть проекта на Flutter Web с поддержкой аутентификации, CRUD-операций, связей \textit{задача-клиент}, базовой аналитики и конфигурационных режимов выполнения.
  \item Проверить корректность работы решения с использованием встроенных сценариев качества (analyze, unit/widget/integration тесты, smoke/E2E-потоки) и сформулировать итоговые выводы о достигнутых результатах.
\end{enumerate}

Объектом исследования в курсовом проекте выступают процессы цифрового управления задачами и клиентскими взаимодействиями в образовательном и малом бизнес-контексте. Предмет исследования --- методы и средства разработки кроссплатформенного CRM-приложения на базе Flutter, Riverpod, GoRouter и Supabase с учетом эксплуатационных требований к надежности, расширяемости и удобству использования \cite{supabaseDocs,postgresDocs}.

Практическая значимость проекта заключается в том, что реализованное приложение может применяться как учебно-прикладной стенд для демонстрации полного цикла разработки: от постановки требований до тестирования и подготовки к деплою. Кроме того, сформированная архитектура допускает масштабирование в сторону более сложных бизнес-сценариев: сегментации клиентов, автоматизации воронки продаж, расширенной отчетности, уведомлений и интеграции внешних каналов взаимодействия \cite{crmBook,projectManagementBook}.

Методологическую основу работы составляют: системный анализ требований, модульный принцип проектирования, сравнительный анализ технологических подходов, итеративная реализация и верификация качества через тестовые сценарии. При подготовке курсового проекта использовались официальные технические руководства по Flutter, Dart, Riverpod, GoRouter и Supabase, а также учебные и научно-практические источники по проектированию информационных систем и программной архитектуре \cite{flutterDocs,dartDocs,riverpodDocs,goRouterDocs,supabaseDocs,isDesign2025}.

Ключевым организационно-техническим принципом проекта является разделение рабочих режимов выполнения. В базовом сценарии используется backend-first подход с реальной авторизацией и операциями чтения/записи в Supabase. Для учебной демонстрации также предусмотрен локальный demo-режим с подстановкой тестовых данных, что позволяет воспроизводить функционал даже при отсутствии внешнего окружения. Такое решение повышает устойчивость проекта в условиях лабораторной и защитной демонстрации, но при этом требует явного разграничения демо и production-потоков для минимизации риска смешения данных \cite{repoReadme,owaspM4}.

Отдельное внимание в работе уделяется качеству программного продукта. Проверка корректности не ограничивается ручным тестированием пользовательского интерфейса. В проект включены команды статического анализа, модульные и виджет-тесты, интеграционные smoke-сценарии и web E2E-проверка основного бизнес-потока \textit{Login $\rightarrow$ Create Client $\rightarrow$ Create Task $\rightarrow$ Verify bindings}. Такой набор позволяет обнаруживать ошибки на разных уровнях системы и снижает вероятность регрессий при дальнейшем развитии кода \cite{repoReadme,testingBook}.

С точки зрения структуры, курсовой проект состоит из введения, двух разделов основной части, заключения, списка использованной литературы и приложений. В первом разделе рассматриваются теоретические основы разработки CRM-ориентированных приложений и обосновывается выбор технологического стека. Во втором разделе представлена практическая реализация \textit{Gamified Task CRM}: архитектура, модель данных, ключевые пользовательские сценарии и результаты верификации. В заключении подведены итоги, дана оценка степени достижения цели и обозначены направления дальнейшего развития системы.

В рамках ограничений курсового проекта акцент сделан на базовом эксплуатационном контуре: корректная работа аутентификации, управление задачами, учет клиентов, формирование отчетных представлений и воспроизводимость CI-проверок. Расширенные возможности (глубокая аналитика, многоканальные уведомления, продвинутая сегментация) рассматриваются как логические точки роста и не противоречат текущей архитектуре.

Таким образом, выбранная тема и реализуемый программный продукт соответствуют требованиям учебной дисциплины, обладают практической применимостью и позволяют в полной мере продемонстрировать навыки системного проектирования и разработки программного обеспечения на современном технологическом стеке.
